% §6.4 Counterfactual rewriting intervention, 12-city panel with NDE/TE and
% an implied NIE = TE - NDE mediation reading. Every number is re-pulled or
% recomputed fresh from the canonical n=200 rerun (Portland now clean):
%   results/counterfactual/counterfactual_12city_table.csv
%   results/counterfactual/counterfactual_12city_pooled.json
%   results/counterfactual/{city}.json
% Cross-generator concordance against results/counterfactual_llama70b/{city}.json
% is recomputed here (Spearman, Pearson, sign agreement over the 12 markets),
% not carried over from session memory. v2 reported a fabricated cross-LLM
% concordance (Spearman rho=0.993, sign 12/12, DML rho=1.0) and fabricated
% validation flip rates; all of that is corrected below.

\subsection{Counterfactual rewriting intervention}
\label{sec:counterfactual-12}

The third identification strategy we deploy is a counterfactual rewriting
intervention. We treat the listing description as a realized draw from a
population of plausible descriptions for the underlying property, and we
probe the causal role of language by generating counterfactual descriptions
in which the listing's submarket-targeted language is altered while the
property's structured attributes are held fixed. The intervention is executed
at the language level by a frozen large language model that receives the
original description, a structured attribute summary, and a target submarket,
and returns a rewritten description. We then re-price the rewritten listing
under the same automated valuation model and read two estimands from the
difference in predicted log price. The style-stripped natural direct effect
holds the model's predicted location at its original value and isolates the
direct contribution of the listing's stylistic and semantic content. The
style-swap total effect rewrites the listing toward a held-out submarket and
captures the full change in predicted price, both the direct change in
language and the change the model infers about location. The difference
between the total effect and the natural direct effect is an implied natural
indirect effect that propagates through the location-prediction channel.

\paragraph{Rewriting validation.} Before we report estimates, we verify that
the rewriting procedure executes the intended intervention. Structured
attribute slots, comprising bedrooms, bathrooms, square footage, year built,
lot size, and three derived flags, are preserved in $99.8\%$ of rewrites on
average across the twelve markets, with the lowest preservation rate at
$99.4\%$ in Phoenix and exact preservation in Boston, Washington, Denver, and
Dallas. The submarket classifier flips its prediction away from the original
submarket in $44.7\%$ of rewrites on average, ranging from $18.4\%$ in San
Francisco to $63.5\%$ in New York. These flip rates are considerably higher
than the figures reported in the previous draft, and they reflect the actual
behavior of the classifier on the rewritten text rather than the much smaller
rates that draft asserted. We read the high average flip rate as evidence that
the rewriting does move the listing's targeted language across submarket
boundaries in a large share of cases, which is the precondition for the
location-mediated channel we examine below, and we read the near-perfect slot
preservation as evidence that the structured attributes survive the rewrite.

\paragraph{Per-market natural direct effect.}
Table~\ref{tab:counterfactual-12} reports the per-market natural direct effect
for the twelve metropolitan markets at $n=200$ listings each. The effect is
positive and excludes zero at the $95\%$ level in nine of the twelve markets,
and it reaches economically meaningful magnitude in seven of them. The largest
direct effects are in Atlanta at $+0.189$ on log price, an implied price uplift
of $+20.8\%$, and Boston at $+0.184$ ($+20.3\%$), followed by Denver at
$+0.169$ ($+18.4\%$), New York at $+0.129$ ($+13.7\%$), Dallas at $+0.118$
($+12.5\%$), Portland at $+0.094$ ($+9.8\%$), and Seattle at $+0.052$
($+5.4\%$). San Francisco carries a small but significant positive effect of
$+0.029$, and Phoenix a positive effect that rounds to zero at $+0.0006$.
Washington returns a small negative direct effect of $-0.006$ that excludes
zero, while Philadelphia and Chicago return point estimates near zero with
confidence intervals that straddle it. The pooled random-effects natural direct
effect across the twelve markets is $+0.083$ on log price, with standard error
$0.010$ and a $95\%$ confidence interval of $[+0.063, +0.102]$, so the pooled
direct effect is positive and significant. The between-market heterogeneity is
substantial, with $Q(11) = 2{,}351$ and $I^2 = 99.5\%$, so we report the
random-effects estimate in preference to the fixed-effects estimate
($+0.0007$, SE $0.00003$) and we caution the reader that the markets differ
enough that the pooled figure is a summary of a dispersed distribution rather
than a single generalizable mean.

\paragraph{Per-market total effect.} The style-swap total effect is small in
absolute magnitude in every market and splits in sign across the panel. It is
strictly negative and excludes zero in five markets, strictly positive and
excludes zero in four, and indistinguishable from zero in the remaining three.
The largest negative total effect is in New York at $-0.049$ on log price, an
implied $-4.8\%$, followed by Dallas at $-0.022$, Boston at $-0.011$, Seattle
at $-0.009$, and a negligible $-0.0002$ in Phoenix. On the positive side,
Denver and Philadelphia both reach about $+0.046$ ($+4.7\%$), with San
Francisco at $+0.012$ and Washington at a small but significant $+0.002$.
Atlanta, Chicago, and Portland return total effects whose confidence intervals
include zero. Because the market-level total effects point in opposite
directions and are individually small, they largely cancel in aggregation. The
pooled random-effects total effect is $-0.001$ on log price, with standard
error $0.002$ and a $95\%$ confidence interval of $[-0.005, +0.003]$ that
includes zero, so the pooled total effect is statistically indistinguishable
from zero even though several markets carry significant effects of either sign.

\paragraph{Mediation decomposition.} The pooled natural direct effect of
$+0.083$ and the pooled total effect of $-0.001$ imply a pooled natural
indirect effect through the location-prediction channel of $-0.084$ on log
price. Substantively this says that the listing description carries a direct
semantic signal that the automated valuation model translates into a positive
price contribution of roughly $+8\%$ relative to the style-stripped
counterfactual, but the style swap also moves the model's inferred location for
the listing, and that shift contributes an indirect effect of comparable
magnitude in the opposite direction. The near-zero pooled total effect is the
net of these two channels of similar size and opposite sign. The substantive
direction of this finding matters for the policy discussion in
Section~\ref{sec:conclusion}. In the automated-valuation deployment we analyze,
listing language influences predicted price along two routes that nearly offset
on average, a direct stylistic route that lifts predicted price and a
location-mediated route that lowers it, and the appearance of a small net
effect masks a direct semantic component that is itself positive and
significant in about half of the markets we study. We treat the decomposition
as a mechanism probe rather than as a precise structural estimate, and we read
it as consistent with the dual role of listing text documented by
\citet{shen2021information} and \citet{baur2023automated}, in which descriptive
language carries both an own-effect on valuation and information that proxies
for location.

\paragraph{Cross-generator robustness.} We re-ran the rewriting intervention
with a second open-weight generator, Llama 3.3 $70$B Instruct, to gauge how far
the conclusions depend on the choice of rewriting model, and we compare it
against the primary Qwen 2.5 $32$B Instruct run. The two runs are not a matched
replication. The Llama run processes $n=100$ listings per market against the
$n=200$ of the primary run, so the comparison is an unmatched cross-generator
probe of direction and rough magnitude rather than a tight correlation between
estimates computed on identical samples. Across the twelve markets the
per-market natural direct effects from the two generators agree in sign in nine
of the twelve markets, with a Spearman rank correlation of $0.55$ and a Pearson
correlation of $0.25$ on the point estimates. The per-market total effects
agree in sign in six of the twelve markets, with a Spearman correlation of
$0.57$ and a Pearson correlation of $0.72$. The central tendencies move in the
same direction across generators for the natural direct effect, with a
positive average in both runs, the random-effects pooled $+0.083$ under Qwen
and a per-market mean of $+0.070$ under Llama. The total effect is the less
stable of the two estimands across generators, near zero under Qwen and
modestly negative under Llama, with a per-market mean of $-0.021$, which is
unsurprising given that the total effect is the small residual after the direct
and indirect channels nearly cancel and is therefore sensitive to the precise
rewriting behavior of each generator. We read the cross-generator evidence as
directional support for the positive natural direct effect and for a near-zero
to slightly negative total effect, and we are explicit that the rank and linear
correlations are moderate rather than near-perfect, so we make no claim that the
two generators reproduce each other's market-level estimates to high precision.
